% ========== PERFORMANCE TARGETS ==========
% Latency requirements, throughput targets, resource budgets, scalability goals
% Source: Symphony/Content/Architecture documents, Performance benchmarking

\section*{Performance Targets}
\addcontentsline{toc}{section}{Performance Targets}

\lettrine{P}{erformance} is a critical differentiator for Symphony, enabling developers to maintain flow state through responsive interactions while supporting enterprise-scale workloads. The performance targets are designed to deliver sub-second response times for interactive operations while supporting thousands of concurrent AI model executions and workflow orchestrations.

\subsection*{Latency Requirements}
\addcontentsline{toc}{subsection}{Latency Requirements}

Symphony's latency requirements are categorized by operation type and user interaction patterns, ensuring optimal responsiveness across all usage scenarios.

\subsubsection*{Interactive Operation Latency}

Operations that directly impact user experience require the most stringent latency requirements to maintain developer productivity and flow state.

\begin{table}[h]
\centering
\caption{Interactive Operation Latency Targets}
\label{tab:interactive-latency}
\begin{tabular}{@{}llllll@{}}
\toprule
\textbf{Operation} & \textbf{Target} & \textbf{P95} & \textbf{P99} & \textbf{Max} & \textbf{Timeout} \\
\midrule
Keystroke Response & 16ms & 25ms & 50ms & 100ms & 200ms \\
Auto-completion & 50ms & 100ms & 200ms & 500ms & 1s \\
Syntax Highlighting & 10ms & 20ms & 50ms & 100ms & 200ms \\
File Operations & 100ms & 200ms & 500ms & 1s & 5s \\
Extension Loading & 200ms & 500ms & 1s & 2s & 10s \\
\bottomrule
\end{tabular}
\end{table}

\textbf{Latency Optimization Strategies}:
\begin{compactlist}
    \item \textbf{Predictive Caching}: Pre-load frequently accessed resources based on usage patterns
    \item \textbf{Incremental Processing}: Break large operations into smaller, interruptible chunks
    \item \textbf{Background Preparation}: Prepare resources in background threads before they're needed
    \item \textbf{Priority Queuing}: Prioritize interactive operations over background tasks
\end{compactlist}

\subsubsection*{AI Model Inference Latency}

AI model operations have varying latency requirements based on their integration into the development workflow.

\begin{table}[h]
\centering
\caption{AI Model Inference Latency Targets}
\label{tab:ai-latency}
\begin{tabular}{@{}llllll@{}}
\toprule
\textbf{Model Type} & \textbf{Target} & \textbf{P95} & \textbf{P99} & \textbf{Max} & \textbf{Use Case} \\
\midrule
Auto-completion & 50ms & 100ms & 200ms & 500ms & Real-time suggestions \\
Code Analysis & 200ms & 500ms & 1s & 2s & Background analysis \\
Code Generation & 1s & 2s & 5s & 10s & Interactive generation \\
Documentation & 2s & 5s & 10s & 30s & Background processing \\
Refactoring & 5s & 10s & 30s & 60s & Complex transformations \\
\bottomrule
\end{tabular}
\end{table}

\textbf{AI Performance Optimization}:
\begin{compactlist}
    \item \textbf{Model Pre-warming}: Keep frequently used models loaded and ready
    \item \textbf{Batch Processing}: Combine multiple requests for improved throughput
    \item \textbf{Result Caching}: Cache model results for identical or similar inputs
    \item \textbf{Progressive Enhancement}: Provide quick results first, then enhance with slower models
\end{compactlist}

\subsubsection*{System Operation Latency}

Core system operations that support the overall architecture have specific latency requirements to ensure system responsiveness.

\begin{table}[h]
\centering
\caption{System Operation Latency Targets}
\label{tab:system-latency}
\begin{tabular}{@{}llllll@{}}
\toprule
\textbf{Operation} & \textbf{Target} & \textbf{P95} & \textbf{P99} & \textbf{Max} & \textbf{Impact} \\
\midrule
IPC Message & 100ns & 200ns & 500ns & 1ms & High \\
Extension Call & 1ms & 2ms & 5ms & 10ms & High \\
Database Query & 1ms & 5ms & 10ms & 50ms & Medium \\
Network Request & 10ms & 50ms & 100ms & 500ms & Medium \\
Disk I/O & 1ms & 10ms & 50ms & 200ms & Low \\
\bottomrule
\end{tabular}
\end{table}

\subsection*{Throughput Targets}
\addcontentsline{toc}{subsection}{Throughput Targets}

Symphony's throughput targets ensure the system can handle enterprise-scale workloads while maintaining responsive performance for individual users.

\subsubsection*{Concurrent User Support}

The system is designed to support large numbers of concurrent users with consistent performance characteristics.

\textbf{User Concurrency Targets}:
\begin{compactlist}
    \item \textbf{Single Instance}: 1,000 concurrent users with full functionality
    \item \textbf{Clustered Deployment}: 10,000+ concurrent users with horizontal scaling
    \item \textbf{Enterprise Deployment}: 100,000+ users with distributed architecture
    \item \textbf{Cloud Deployment}: Unlimited scaling with auto-scaling capabilities
\end{compactlist}

\textbf{Per-User Resource Allocation}:
\begin{compactlist}
    \item \textbf{Memory}: 50-150MB per active user session
    \item \textbf{CPU}: 0.1-0.5 CPU cores per active user (burst to 2 cores)
    \item \textbf{Storage}: 100MB-1GB per user workspace
    \item \textbf{Network}: 1-10 Mbps per user for real-time operations
\end{compactlist}

\subsubsection*{AI Model Throughput}

AI model execution throughput is critical for supporting multiple concurrent development workflows.

\begin{table}[h]
\centering
\caption{AI Model Throughput Targets}
\label{tab:ai-throughput}
\begin{tabular}{@{}llllll@{}}
\toprule
\textbf{Model Category} & \textbf{Requests/sec} & \textbf{Batch Size} & \textbf{GPU Util} & \textbf{Memory} & \textbf{Scaling} \\
\midrule
Small Models (<1B) & 1000 & 32 & 60-80\% & 2-4GB & Linear \\
Medium Models (1-10B) & 100 & 16 & 70-90\% & 8-16GB & Sub-linear \\
Large Models (10B+) & 10 & 4 & 80-95\% & 24-48GB & Logarithmic \\
Specialized Models & 50 & 8 & 65-85\% & 4-12GB & Variable \\
\bottomrule
\end{tabular}
\end{table}

\textbf{Throughput Optimization Techniques}:
\begin{compactlist}
    \item \textbf{Dynamic Batching}: Automatically batch requests to maximize GPU utilization
    \item \textbf{Model Parallelism}: Distribute large models across multiple GPUs
    \item \textbf{Pipeline Parallelism}: Pipeline model execution stages for improved throughput
    \item \textbf{Quantization}: Use quantized models for improved throughput with acceptable quality loss
\end{compactlist}

\subsubsection*{Data Processing Throughput}

High-throughput data processing capabilities support large-scale code analysis and transformation operations.

\textbf{Data Processing Targets}:
\begin{compactlist}
    \item \textbf{File Processing}: 10,000 files/second for syntax analysis
    \item \textbf{Code Parsing}: 1 million lines/second for AST generation
    \item \textbf{Index Updates}: 100,000 symbols/second for code indexing
    \item \textbf{Search Operations}: 10,000 queries/second with sub-second response
\end{compactlist}

\subsection*{Resource Budgets}
\addcontentsline{toc}{subsection}{Resource Budgets}

Symphony's resource budgets ensure efficient utilization of system resources while maintaining performance and reliability guarantees.

\subsubsection*{Memory Budget Allocation}

Memory allocation is carefully managed to prevent resource exhaustion while supporting concurrent operations.

\begin{table}[h]
\centering
\caption{Memory Budget Allocation}
\label{tab:memory-budget}
\begin{tabular}{@{}llllll@{}}
\toprule
\textbf{Component} & \textbf{Base} & \textbf{Per User} & \textbf{Peak} & \textbf{Limit} & \textbf{Priority} \\
\midrule
Core System & 500MB & - & 1GB & 2GB & Critical \\
AI Models & 2GB & 50MB & 16GB & 32GB & High \\
Extensions & 100MB & 25MB & 2GB & 4GB & Medium \\
User Data & 50MB & 100MB & 1GB & 2GB & Medium \\
Caches & 200MB & 10MB & 4GB & 8GB & Low \\
\bottomrule
\end{tabular}
\end{table}

\textbf{Memory Management Strategies}:
\begin{compactlist}
    \item \textbf{Lazy Loading}: Load resources only when needed to minimize memory footprint
    \item \textbf{Memory Pooling}: Reuse memory allocations to reduce fragmentation
    \item \textbf{Garbage Collection}: Proactive cleanup of unused resources
    \item \textbf{Memory Compression}: Compress inactive data to reduce memory usage
\end{compactlist}

\subsubsection*{CPU Budget Allocation}

CPU resources are allocated based on operation priority and user interaction requirements.

\textbf{CPU Allocation Strategy}:
\begin{compactlist}
    \item \textbf{Interactive Operations}: 40\% of CPU reserved for user interactions
    \item \textbf{AI Model Execution}: 35\% allocated for AI model processing
    \item \textbf{Background Tasks}: 20\% for indexing, analysis, and maintenance
    \item \textbf{System Operations}: 5\% reserved for system management
\end{compactlist}

\textbf{CPU Optimization Techniques}:
\begin{compactlist}
    \item \textbf{Work Stealing}: Idle threads steal work from busy threads
    \item \textbf{NUMA Awareness}: Schedule tasks on appropriate NUMA nodes
    \item \textbf{CPU Affinity}: Pin performance-critical threads to specific cores
    \item \textbf{Dynamic Scaling}: Adjust thread pools based on current load
\end{compactlist}

\subsubsection*{Storage Budget Management}

Storage resources are managed to balance performance, capacity, and cost considerations.

\begin{infobox}[title=Storage Tier Strategy]
\textbf{Hot Storage (SSD)}: Frequently accessed data, active projects, and AI model cache

\textbf{Warm Storage (Fast HDD)}: Recent projects, extension cache, and temporary files

\textbf{Cold Storage (Archive)}: Historical data, backups, and rarely accessed content

\textbf{Cloud Storage}: Overflow capacity, collaboration data, and disaster recovery
\end{infobox}

\subsection*{Scalability Goals}
\addcontentsline{toc}{subsection}{Scalability Goals}

Symphony's scalability goals ensure the system can grow to meet increasing demands while maintaining performance and reliability characteristics.

\subsubsection*{Horizontal Scaling Capabilities}

The architecture supports horizontal scaling across multiple dimensions to handle increased load.

\textbf{Scaling Dimensions}:
\begin{expandedlist}
    \item \textbf{User Scaling}: Linear scaling of user capacity through load balancing and session distribution
    \item \textbf{Compute Scaling}: Elastic scaling of AI model execution through GPU cluster management
    \item \textbf{Storage Scaling}: Distributed storage scaling through sharding and replication
    \item \textbf{Network Scaling}: CDN integration and edge computing for global performance
\end{expandedlist}

\textbf{Auto-Scaling Triggers}:
\begin{compactlist}
    \item \textbf{CPU Utilization}: Scale when average CPU usage exceeds 70\% for 5 minutes
    \item \textbf{Memory Pressure}: Scale when memory usage exceeds 80\% of available capacity
    \item \textbf{Queue Depth}: Scale when task queues exceed 100 pending items
    \item \textbf{Response Time}: Scale when P95 response time exceeds target by 50\%
\end{compactlist}

\subsubsection*{Performance Scaling Characteristics}

Understanding how performance scales with system size is crucial for capacity planning and architecture decisions.

\begin{table}[h]
\centering
\caption{Performance Scaling Characteristics}
\label{tab:scaling-performance}
\begin{tabular}{@{}llllll@{}}
\toprule
\textbf{Metric} & \textbf{1x} & \textbf{10x} & \textbf{100x} & \textbf{1000x} & \textbf{Scaling} \\
\midrule
User Capacity & 1K & 10K & 100K & 1M & Linear \\
AI Throughput & 100/s & 800/s & 6K/s & 40K/s & Sub-linear \\
Storage IOPS & 10K & 80K & 600K & 4M & Sub-linear \\
Network Bandwidth & 1Gbps & 8Gbps & 60Gbps & 400Gbps & Sub-linear \\
\bottomrule
\end{tabular}
\end{table}

\subsubsection*{Enterprise Scalability Requirements}

Enterprise deployments require specific scalability characteristics to support large organizations.

\textbf{Enterprise Scale Targets}:
\begin{compactlist}
    \item \textbf{Organization Size}: Support for 100,000+ developers in single deployment
    \item \textbf{Project Scale}: Handle projects with millions of files and billions of lines of code
    \item \textbf{Geographic Distribution}: Multi-region deployment with <100ms cross-region latency
    \item \textbf{High Availability}: 99.99\% uptime with automatic failover and disaster recovery
\end{compactlist}

\begin{successbox}
\textbf{Performance Monitoring}: Symphony includes comprehensive performance monitoring and alerting to ensure performance targets are met in production. Real-time dashboards provide visibility into all key performance metrics, enabling proactive optimization and capacity planning.
\end{successbox}

\subsubsection*{Performance Benchmarking}

Regular performance benchmarking ensures Symphony meets its performance targets across different deployment scenarios.

\textbf{Benchmark Categories}:
\begin{compactlist}
    \item \textbf{Synthetic Benchmarks}: Controlled tests of specific system components
    \item \textbf{Real-World Workloads}: Tests based on actual user usage patterns
    \item \textbf{Stress Testing}: Maximum capacity testing under extreme load conditions
    \item \textbf{Regression Testing}: Continuous testing to prevent performance degradation
\end{compactlist}

\textbf{Performance Validation Process}:
\begin{compactlist}
    \item \textbf{Continuous Integration}: Automated performance tests on every code change
    \item \textbf{Staging Environment}: Full-scale performance testing before production deployment
    \item \textbf{Production Monitoring}: Continuous monitoring of production performance metrics
    \item \textbf{Performance Reviews}: Regular review of performance trends and optimization opportunities
\end{compactlist}

\begin{alertbox}
\textbf{Performance Trade-offs}: While Symphony's performance targets are ambitious, they represent a balance between performance, functionality, and resource efficiency. Some advanced AI features may require relaxed performance targets in exchange for enhanced capabilities.
\end{alertbox}